\documentclass[journal]{IEEEtran}
\usepackage{amsmath,amssymb,amsthm,booktabs,graphicx,etoolbox,url,hyperref}
\IfFileExists{generated/numbers.tex}{\newcommand{\TotalRuns}{640}
\newcommand{\diabetesmlpFedavgAuroc}{0.6203}
\newcommand{\diabetesmlpFedavgAuprc}{0.171}
\newcommand{\diabetesmlpFedavgUploadMb}{218.4}
\newcommand{\diabetesmlpFedavgSeeds}{10}
\newcommand{\diabetesmlpFedavgVsFedperDiff}{-0.0012}
\newcommand{\diabetesmlpFedavgNiAurocP}{$<10^{-4}$}
\newcommand{\diabetesmlpFedavgNiAuprcP}{$<10^{-4}$}
\newcommand{\diabetesmlpFedavgUploadSavingPct}{-0}
\newcommand{\diabetesmlpFedproxAuroc}{0.6279}
\newcommand{\diabetesmlpFedproxAuprc}{0.175}
\newcommand{\diabetesmlpFedproxUploadMb}{218.4}
\newcommand{\diabetesmlpFedproxSeeds}{10}
\newcommand{\diabetesmlpFedproxVsFedperDiff}{+0.0065}
\newcommand{\diabetesmlpFedproxNiAurocP}{$<10^{-4}$}
\newcommand{\diabetesmlpFedproxNiAuprcP}{$<10^{-4}$}
\newcommand{\diabetesmlpFedproxUploadSavingPct}{-0}
\newcommand{\diabetesmlpScaffoldAuroc}{0.5930}
\newcommand{\diabetesmlpScaffoldAuprc}{0.138}
\newcommand{\diabetesmlpScaffoldUploadMb}{436.9}
\newcommand{\diabetesmlpScaffoldSeeds}{10}
\newcommand{\diabetesmlpScaffoldVsFedperDiff}{-0.0284}
\newcommand{\diabetesmlpScaffoldNiAurocP}{0.997}
\newcommand{\diabetesmlpScaffoldNiAuprcP}{1.000}
\newcommand{\diabetesmlpScaffoldUploadSavingPct}{-101}
\newcommand{\diabetesmlpFedperAuroc}{0.6215}
\newcommand{\diabetesmlpFedperAuprc}{0.171}
\newcommand{\diabetesmlpFedperUploadMb}{217.8}
\newcommand{\diabetesmlpFedperSeeds}{10}
\newcommand{\diabetesmlpFedlamaAuroc}{0.6291}
\newcommand{\diabetesmlpFedlamaAuprc}{0.176}
\newcommand{\diabetesmlpFedlamaUploadMb}{154.1}
\newcommand{\diabetesmlpFedlamaSeeds}{10}
\newcommand{\diabetesmlpFedlamaVsFedperDiff}{+0.0076}
\newcommand{\diabetesmlpFedlamaNiAurocP}{$<10^{-4}$}
\newcommand{\diabetesmlpFedlamaNiAuprcP}{$<10^{-4}$}
\newcommand{\diabetesmlpFedlamaUploadSavingPct}{29}
\newcommand{\diabetesmlpAdacommAuroc}{0.6302}
\newcommand{\diabetesmlpAdacommAuprc}{0.168}
\newcommand{\diabetesmlpAdacommUploadMb}{29.1}
\newcommand{\diabetesmlpAdacommSeeds}{10}
\newcommand{\diabetesmlpAdacommVsFedperDiff}{+0.0087}
\newcommand{\diabetesmlpAdacommNiAurocP}{$<10^{-4}$}
\newcommand{\diabetesmlpAdacommNiAuprcP}{0.000}
\newcommand{\diabetesmlpAdacommUploadSavingPct}{87}
\newcommand{\diabetesmlpLagAuroc}{0.6351}
\newcommand{\diabetesmlpLagAuprc}{0.179}
\newcommand{\diabetesmlpLagUploadMb}{218.4}
\newcommand{\diabetesmlpLagSeeds}{10}
\newcommand{\diabetesmlpLagVsFedperDiff}{+0.0136}
\newcommand{\diabetesmlpLagNiAurocP}{$<10^{-4}$}
\newcommand{\diabetesmlpLagNiAuprcP}{$<10^{-4}$}
\newcommand{\diabetesmlpLagUploadSavingPct}{-0}
\newcommand{\diabetesmlpFixedAuroc}{0.6300}
\newcommand{\diabetesmlpFixedAuprc}{0.175}
\newcommand{\diabetesmlpFixedUploadMb}{128.4}
\newcommand{\diabetesmlpFixedSeeds}{10}
\newcommand{\diabetesmlpFixedVsFedperDiff}{+0.0085}
\newcommand{\diabetesmlpFixedNiAurocP}{$<10^{-4}$}
\newcommand{\diabetesmlpFixedNiAuprcP}{$<10^{-4}$}
\newcommand{\diabetesmlpFixedUploadSavingPct}{41}
\newcommand{\diabetesmlpTierfedAuroc}{0.6269}
\newcommand{\diabetesmlpTierfedAuprc}{0.175}
\newcommand{\diabetesmlpTierfedUploadMb}{65.3}
\newcommand{\diabetesmlpTierfedSeeds}{10}
\newcommand{\diabetesmlpTierfedVsFedperDiff}{+0.0054}
\newcommand{\diabetesmlpTierfedNiAurocP}{$<10^{-4}$}
\newcommand{\diabetesmlpTierfedNiAuprcP}{$<10^{-4}$}
\newcommand{\diabetesmlpTierfedUploadSavingPct}{70}
\newcommand{\diabetesmlpNTrain}{55,992}
\newcommand{\diabetesmlpNTest}{13,998}
\newcommand{\diabetesmlpNFeatures}{160}
\newcommand{\diabetesmlpNPatients}{55,992}
\newcommand{\diabetesmlpPrevalence}{9.0}
\newcommand{\diabetesmlpMaxClientShare}{12}
\newcommand{\diabetesmlpLabelTv}{0.009}
\newcommand{\diabetesmlpTierSlowParams}{20,608}
\newcommand{\diabetesmlpTierSlowPct}{45.3}
\newcommand{\diabetesmlpTierMediumParams}{24,768}
\newcommand{\diabetesmlpTierMediumPct}{54.4}
\newcommand{\diabetesmlpTierFastParams}{130}
\newcommand{\diabetesmlpTierFastPct}{0.3}
\newcommand{\diabetesmlpTierTotalParams}{45,506}
\newcommand{\diabetesmlpTierTotalPct}{100.0}
\newcommand{\flagdiabetesmlpFedavgNonInferiorAuroc}{true}
\newcommand{\flagdiabetesmlpFedavgNonInferiorAuprc}{true}
\newcommand{\flagdiabetesmlpFedproxNonInferiorAuroc}{true}
\newcommand{\flagdiabetesmlpFedproxNonInferiorAuprc}{true}
\newcommand{\flagdiabetesmlpScaffoldNonInferiorAuroc}{false}
\newcommand{\flagdiabetesmlpScaffoldNonInferiorAuprc}{false}
\newcommand{\flagdiabetesmlpFedlamaNonInferiorAuroc}{true}
\newcommand{\flagdiabetesmlpFedlamaNonInferiorAuprc}{true}
\newcommand{\flagdiabetesmlpAdacommNonInferiorAuroc}{true}
\newcommand{\flagdiabetesmlpAdacommNonInferiorAuprc}{true}
\newcommand{\flagdiabetesmlpLagNonInferiorAuroc}{true}
\newcommand{\flagdiabetesmlpLagNonInferiorAuprc}{true}
\newcommand{\flagdiabetesmlpFixedNonInferiorAuroc}{true}
\newcommand{\flagdiabetesmlpFixedNonInferiorAuprc}{true}
\newcommand{\flagdiabetesmlpTierfedNonInferiorAuroc}{true}
\newcommand{\flagdiabetesmlpTierfedNonInferiorAuprc}{true}
}{}
\newcommand{\runreq}{\textbf{[RUN REQUIRED]}}
\newcommand{\gen}[1]{\IfFileExists{generated/#1.tex}{\input{generated/#1.tex}}{\runreq}}
\newcommand{\num}[1]{\ifcsdef{#1}{\csname #1\endcsname}{\runreq}}
\newtheorem{proposition}{Proposition}
\newtheorem{lemma}{Lemma}
\begin{document}
\title{Supplementary Material: Communication as the Binding Constraint}
\author{Abraheem Rashid}
\maketitle

\section{S1. Relation to the superseded version}
An earlier manuscript (``Nested Federated Learning'' / ``TierFed'' v2) reported results produced by code with the following defects, found in an internal audit and corrected here: (1) differential privacy was applied inside two strategies only and clipped the aggregate rather than each client's update, so FedAvg/FedPer/FedProx ran without noise and the reported budgets were invalid; the TierFed DP runs crashed on local-only rounds and the table was unsupported; (2) the SCAFFOLD baseline was FedAvg with doubled bit accounting; (3) the scalability table contained no TierFed runs and its values were not those of the logs; (4) the architecture sweep compared different depths across methods; (5) the imaging result used an undisclosed configuration on a binarised 97\%-positive task; (6) the inversion attack targeted the fixed schedule at initialisation and was labelled TierFed; (7) the model described (93 inputs) was not the model trained (162 inputs); (8) repeated runs were not deterministic; (9) no patient-level deduplication was performed and one ``missing specialty'' client held 49\% of the data; (10) the greedy controller reached only powers of two. Every item maps to a code change in the released repository (\texttt{ROADMAP.md}, section 0) and to a test.

\section{S2. Derivation of Proposition 1}
\label{supp:proof}
We follow the cooperative-SGD framework of Wang and Joshi [JMLR 2021, Thm.~1], which bounds the averaged gradient norm of local-update SGD in terms of the per-round deviation of client iterates from the virtual average. Write the shared parameter as $\theta=(\theta_g)_{g}$ with disjoint blocks. For site $k$ and block $g$ let $\theta^{(t)}_{k,g}$ be the iterate after $t$ local steps since block $g$ was last synchronised and $\bar\theta^{(t)}_g$ the (virtual) average across sites.

\begin{lemma}[Blockwise deviation]
Under A1--A3, for $\eta\le c_0/(LK_{\max}\tau_{\max})$ and any $t\le K_g\tau_{\max}$,
\[
\frac1N\sum_k\mathbb E\big\|\theta^{(t)}_{k,g}-\bar\theta^{(t)}_g\big\|^2\ \le\ c\,\eta^2 t^2\,(\Gamma_g+\sigma_g^2)\ \le\ c\,\eta^2K_g^2\tau_{\max}^2(\Gamma_g+\sigma_g^2).
\]
\end{lemma}
\begin{proof}
Identical to Lemma 2 of Wang and Joshi restricted to the coordinates of block $g$: the deviation after $t$ steps is a sum of $t$ stochastic-gradient differences, each with second moment bounded by $\Gamma_g+\sigma_g^2$ by (A2)--(A3), and the smoothness term is absorbed by the step-size condition. The restriction to a block is valid because the gradient of block $g$ is bounded blockwise by the same assumptions.
\end{proof}

The descent-lemma step of the cited theorem bounds $\mathbb E\|\nabla F(\bar\theta_r)\|^2$ by the progress term, the variance term and $L^2$ times the \emph{total} deviation $\frac1N\sum_k\|\theta_k-\bar\theta\|^2$. Because the blocks are disjoint, the total deviation is the sum over $g$ of the blockwise deviations, each bounded by the lemma with its own $K_g$; summing and averaging over rounds gives Proposition~1 of the main text with constants inherited from the cited theorem. We make no claim about the constants beyond their existence. Setting $K_g=1$ for a single block returns the cited bound with $\tau$ local steps.

\textbf{What the controller sees.} By the same lemma the dispersion of the received slices of tier $g$ after $K_g$ rounds is of order $\eta^2K_g^2\tau^2(\Gamma_g+\sigma_g^2)$; normalising by $\sum_j\|\bar\theta_j\|^2$ gives the scale-free $\hat\Gamma_g$ of Eq.~(7) of the main text (the dispersion estimator). The controller therefore minimises $\sum_g K_g^2\hat\Gamma_g$ where $\hat\Gamma_g$ already carries the factor $K_g^2$ of the interval that produced it; it is a feedback loop on a proxy, and its stability is observed (Fig.~4 of the main text), not proved.

\section{S3. Differential-privacy accounting}
Per-client clipping to $C$ makes the $\ell_2$ sensitivity of the sum $\sum_k\tilde\Delta_k$ to one client exactly $C$. The Gaussian mechanism with standard deviation $\sigma C$ satisfies $(\alpha,\alpha/(2\sigma^2))$-RDP for every $\alpha>1$ [Mironov 2017, Prop.~7]. Composition over the $R_{\mathrm{tx}}$ rounds in which something is transmitted gives $(\alpha,R_{\mathrm{tx}}\alpha/(2\sigma^2))$-RDP, and $(\varepsilon,\delta)$-DP follows with $\varepsilon=\min_\alpha\{R_{\mathrm{tx}}\alpha/(2\sigma^2)+\log(1/\delta)/(\alpha-1)\}$ [Mironov 2017, Prop.~3]. Rounds with no transmission release nothing and are not charged. Full participation is used, so no subsampling amplification applies.

\section{S4. Per-seed tables}
Every run's final metrics are in \texttt{results/tables/collected\_all\_raw.csv} in the released repository. Total runs at build time: \num{TotalRuns}.

\section{S5. SGD arm}
Identical schedules with local SGD (\texttt{--optimizer sgd}, learning rate $10^{-2}$) on the main benchmark: \runreq{} (generated when the arm is run).

\section{S6. Excluded sites on the imaging benchmark}
Number of sites whose patient-grouped held-out split contained a single class (identical for every method, per seed): \runreq{} (read from \texttt{n\_clients\_single\_class\_final} in the logs when the imaging stage has run).

\end{document}
